\section{GAN 把散度的估计交给判别器}
\label{sec:14-Deep-Learning/07-Generative-Models-Unified/03}

训练曲线上判别器一路走高：真假分辨准确率 \(99\%\)，判别损失贴着零。按``对手越强越能逼出好选手''的直觉，这应该是好消息。实际发生的是生成器的梯度范数掉了三个数量级，生成的图三万步没有任何肉眼可见的变化。有人把判别器调弱——每轮少更新几步、给输入加一点噪声、把学习率砍半——生成器立刻又开始动了。

一个组件变强让整个系统停摆，说明``判别器是生成器的对手''这个说法在关键处不成立。它不是对手。它是一个测量仪器，而仪器读数变得极端，恰恰说明被测的那个量已经没有信息可给了。

\subsection{最优判别器把目标变成 JS 散度}

原始的极小极大目标是

\[
\min_\theta\ \max_D\ V(D,\theta)
=\E_{x\sim p_{\text{data}}}\!\left[\log D(x)\right]
+\E_{x\sim p_\theta}\!\left[\log\left(1-D(x)\right)\right],
\]

其中 \(D:\R^d\to(0,1)\)，\(p_\theta\) 是生成器把噪声推过去得到的分布。这个式子里没有任何密度出现——两个期望都只需要能从对应分布采样，这正是 GAN 放弃密度换来的自由。

固定生成器，先解内层。把两个期望写成同一个积分：

\[
V(D,\theta)=\int\Big[p_{\text{data}}(x)\log D(x)+p_\theta(x)\log\left(1-D(x)\right)\Big]dx.
\]

被积函数在每个 \(x\) 上只依赖 \(D(x)\) 这一个数。如果 \(D\) 可以在不同的 \(x\) 上独立取值，那么最大化整个积分就等于对每个 \(x\) 分别最大化方括号里的东西。记 \(a=p_{\text{data}}(x)\)、\(b=p_\theta(x)\)，要在 \(t\in(0,1)\) 上最大化 \(h(t)=a\log t+b\log(1-t)\)。求导：

\[
h'(t)=\frac{a}{t}-\frac{b}{1-t}=0
\quad\Longrightarrow\quad
t=\frac{a}{a+b},
\]

而 \(h''(t)=-a/t^2-b/(1-t)^2<0\)，所以这个驻点是唯一的极大点。于是

\[
D^\ast(x)=\frac{p_{\text{data}}(x)}{p_{\text{data}}(x)+p_\theta(x)}.
\]

代回去。记混合分布 \(m=\tfrac12(p_{\text{data}}+p_\theta)\)：

\[
\begin{aligned}
V(D^\ast,\theta)
&=\int p_{\text{data}}\log\frac{p_{\text{data}}}{p_{\text{data}}+p_\theta}
+p_\theta\log\frac{p_\theta}{p_{\text{data}}+p_\theta}\\
&=\int p_{\text{data}}\log\frac{p_{\text{data}}}{2m}
+p_\theta\log\frac{p_\theta}{2m}\\
&=\KL(p_{\text{data}}\,\Vert\,m)+\KL(p_\theta\,\Vert\,m)-\log 4\\
&=2\operatorname{JS}(p_{\text{data}}\,\Vert\,p_\theta)-\log 4 .
\end{aligned}
\]

条件对账。逐点优化那一步要求 \(D\) 所在的函数类足够大，能在每个 \(x\) 上独立取到 \(a/(a+b)\)——有限容量的网络并不满足，所以实际训练里的判别器给出的只是 \(D^\ast\) 的一个近似，代回后得到的也只是 JS 的一个下界估计。求导那一步要求 \(a+b>0\)，即该点至少有一个分布放了质量；\(a=b=0\) 的地方被积函数恒为零，\(D\) 取什么都无所谓。第三行用到 \(\int p_{\text{data}}=\int p_\theta=1\) 把 \(\log 2\) 提出来。

结论只有一句：生成器在最小化 JS 散度，而判别器是那个散度的估计器。你算不出 \(p_\theta\)，于是训一个网络把 \(D^\ast\) 逼出来，再从它身上把散度的值和梯度取回来。JS 的性质在 \hyperref[sec:07-Information-Theory/03-Divergences/03]{``Jensen--Shannon 散度与尺度选择''} 里已经建立，这里要用的是它的一个具体缺陷。

\subsection{支撑不重叠时，JS 没有梯度可给}

生成器把一个 \(k\) 维噪声推成 \(d\) 维样本，\(k\) 通常远小于 \(d\)。上一篇算过这一步：像集是 \(\R^d\) 里一个至多 \(k\) 维的集合，Lebesgue 测度为零。数据本身按流形假设也集中在一个低维集合附近。两个低维集合在高维空间里处于一般位置时，交集的测度同样是零。

支撑不重叠意味着：对几乎所有 \(x\)，\(p_{\text{data}}(x)\) 与 \(p_\theta(x)\) 至多一个非零。代进 JS 的定义，每一项都取到 \(\log 2\)，于是

\[
\operatorname{JS}(p_{\text{data}}\,\Vert\,p_\theta)=\log 2
\]

——一个常数，与参数无关。常数的梯度是零。同时 \(D^\ast\) 在数据侧取 \(1\)、在生成侧取 \(0\)，判别器完美，损失贴零。开头那张训练曲线上的两件事，是同一件事的两个读数。

最小的例子把这一点摆得再清楚不过。取 \(p_0\) 是集中在原点的单点分布，\(p_\vartheta\) 集中在 \((\vartheta,0)\)。只要 \(\vartheta\neq0\)，两个分布支撑不交，\(\operatorname{JS}=\log 2\)；\(\vartheta=0\) 时突然变成 \(0\)。这个函数在 \(\vartheta=0\) 处不连续，在别处导数为零，任何基于梯度的方法都无从知道该往哪边挪。而搬运质量的代价 \(\abs{\vartheta}\) 在同一情形下处处给出正确的方向。

\begin{center}
\begin{tikzpicture}[>=stealth]
\begin{scope}[x=1.02cm,y=1.05cm]
  \draw[->] (-0.15,0) -- (3.35,0) node[right] {\(x\)};
  \draw[blue!70!black,very thick,->] (1,0) -- (1,1.15);
  \draw[red!75!black,very thick,->] (2,0) -- (2,1.15);
  \node[blue!70!black,font=\small] at (0.62,1.22) {\(p_{\text{data}}=\delta_1\)};
  \node[red!75!black,font=\small] at (2.42,1.22) {\(p_\theta=\delta_2\)};
  \draw[<->,gray!85] (1,-0.22) -- (2,-0.22);
  \node[gray!60!black,font=\small] at (1.5,-0.5) {间距 \(d>0\)，两个支撑不交};
\end{scope}
\begin{scope}[shift={(4.7,0)},x=1.55cm,y=1.5cm]
  \draw[->] (-0.08,0) -- (1.72,0) node[right] {\(d\)};
  \draw[->] (0,-0.08) -- (0,1.45);
  \draw[gray!55,dashed] (0,0.693) -- (1.62,0.693);
  \node[gray!55!black,font=\small,left] at (0,0.693) {\(\log 2\)};
  \fill[blue!70!black] (0,0) circle (1.7pt);
  \draw[blue!70!black,thick] (0.055,0.6931) -- (1.55,0.6931);
  \draw[blue!70!black,fill=white,thick] (0,0.6931) circle (1.7pt);
  \draw[red!75!black,thick] (0,0) -- (1.2,1.2);
  \node[blue!70!black,font=\small,below] at (1.15,0.66) {JS};
  \node[red!75!black,font=\small,above left] at (1.22,1.20) {Wasserstein};
\end{scope}
\end{tikzpicture}

\emph{对图中的两个点质量，\(d=0\) 时 JS 为 \(0\)，任意 \(d>0\) 时支撑不交，JS 立刻跳到常数 \(\log2\)；搬运质量的代价却继续按 \(d\) 线性增长。梯度有没有，看的是曲线平不平，不是散度大不大。}
\end{center}

图上那条平线解释了为什么``把判别器调弱''会有效：容量不足或加了噪声的判别器估不出那个尖锐的 \(D^\ast\)，实际给出的是一个被抹平的版本，对应的散度估计因此还留着一点斜率。这是在用估计误差换梯度信号，不是在改进目标。

\subsection{换一把尺子：搬运代价而不是重叠程度}

\(1\)-Wasserstein 距离量的是另一件事：

\[
W_1(p,q)=\inf_{\gamma\in\Pi(p,q)}\ \E_{(x,y)\sim\gamma}\!\left[\norm{x-y}\right],
\]

其中 \(\Pi(p,q)\) 是所有边缘分别为 \(p\) 和 \(q\) 的联合分布。它问的是把 \(p\) 这堆土搬成 \(q\) 那堆土最少要走多少路，而不问两堆土重叠了多少。支撑完全不交时它仍然是个有限的、随位置平滑变化的数——上面那个单点例子里恰好等于 \(\abs{\vartheta}\)。

但这个下确界是在一个无穷维的联合分布集合上取的，直接算不了。Kantorovich--Rubinstein 对偶把它翻到另一侧：

\[
W_1(p,q)=\sup_{f:\ \operatorname{Lip}(f)\le1}\ \E_{x\sim p}\!\left[f(x)\right]-\E_{x\sim q}\!\left[f(x)\right].
\]

右边只用到两个分布下的期望——又是只需要采样。于是判别器改了岗位：它不再输出一个概率，而是输出一个实数打分，目标从交叉熵换成两侧打分的差；对应地它也不再叫判别器，叫 critic。

换来这一切的是那个约束：\(f\) 必须 \(1\)-Lipschitz，否则上确界发散，右边等于 \(+\infty\)。落地的手段有三种，每种都只是近似满足这个条件。权重裁剪把每个参数硬夹在一个小区间里，确实保证了有限的 Lipschitz 常数，但同时把函数族削得很窄，critic 往往退化成接近分片线性的简单函数。梯度惩罚在数据点与生成点之间的随机插值上加一项 \((\norm{\nabla_x f(x)}-1)^2\)，只在采样到的那些点上生效，是软约束，别处不管。谱归一化把每层权重矩阵除以它的最大奇异值，使每层的 Lipschitz 常数不超过 \(1\)，整网的常数因而不超过各层之积——最大奇异值就是矩阵的算子范数，它作为``最大拉伸倍数''的意义在 \hyperref[sec:03-Linear-Algebra/07-SVD-and-Data-Applications/01]{``任意矩阵的奇异值分解''} 里说清楚了。这个界通常很松：各层最坏方向不会同时对齐，实际 Lipschitz 常数往往远小于上界。

所以``WGAN 在优化 Wasserstein 距离''是一句近似陈述。约束满足到什么程度，对偶就成立到什么程度，报告里给出的那个 critic 差值不能直接当成 \(W_1\) 的估计值来跨模型比较。

\subsection{模式坍缩不是新现象}

第 01 篇那条只贴住一个峰的红色曲线，在这里落地。

名义上 GAN 优化的 JS 是对称的，两个方向的惩罚都该在。可实际训练有两处把它推向反向 KL 那一侧。

第一处是损失的改写。原始生成器目标 \(\min_\theta\E_{p_\theta}[\log(1-D(x))]\) 在训练早期判别器占优时梯度接近零（\(D\to0\) 时 \(\log(1-D)\) 几乎是平的），于是实践中一律改成最大化 \(\E_{p_\theta}[\log D(x)]\)。这不是同一个目标的等价变形：在最优判别器下，这个非饱和目标对应的量里含有一个 \(-\KL(p_\theta\Vert p_{\text{data}})\) 项，也就是反向 KL 的方向，天然是模式寻求的。

第二处更本质。判别器只对\emph{采到的样本}打分。生成器输出的东西不像真的，它会给低分；某一类真实样本从来没被生成过，它没有任何机制去报告这件事——没生成出来的东西不会出现在它的输入里。目标函数里因此缺少``漏掉整个模式''这一项的惩罚，而第 01 篇说过，不被惩罚的行为就是被允许的行为。

收缩到少数模式，于是不是训练没跑好，是目标函数认可的一个解。加噪声、加多样性正则、用小批量统计量喂给判别器，这些做法之所以有时管用，是因为它们人为地把``漏掉模式''这件事重新写回了目标里。

\subsection{哪些结论是定理，哪些只是经验}

上面所有代数推导都挂在同一个假设上：每一步生成器更新之前，判别器已经达到最优。实际训练里两者交替走小步，得到的是一个非凸非凹的极小极大问题，而这类问题连收敛都不该被默认。

一个两行的反例足以说明性质有多不同。取 \(V(x,y)=xy\)，同时对 \(x\) 做梯度下降、对 \(y\) 做梯度上升，连续时间的动力学是

\[
\dot x=-y,\qquad \dot y=x,
\]

解是绕原点的匀速圆周运动，\(x^2+y^2\) 严格守恒。唯一的鞍点 \((0,0)\) 就在正中间，轨迹永远绕着它转、一步也不靠近。离散化之后半径还会缓慢增大。这里没有任何非凸性、没有深度网络、没有随机性，纯粹是极小极大结构本身的性质。

因此实践中那套稳定训练的做法——判别器与生成器用不同的学习率、对生成器权重做指数滑动平均、给判别器加谱归一化或梯度惩罚、对判别器输入做可微的数据增广——都是经验规律。它们在多个数据集上被反复复现，值得作为默认起点，但没有定理保证它们会让训练收敛，也没有理由认为在一个新的数据分布上它们必然还成立。这一段与前面几节的地位不同，不要混着引用。

\subsection{往下走之前，先把``好''说清楚}

三条路线现在都摆开了：流要精确似然，付出架构自由；VAE 要可训练的下界，付出模糊；GAN 要架构自由和锐利样本，付出理论保证和模式覆盖。它们的取舍方向不同，可判断谁更好之前，得先有一把能同时量到``像不像''和``全不全''的尺子——而单一数字给不出这两样。下一篇处理评价，然后才是那个把一次远距离匹配拆成许多次近距离匹配的想法。
